\subsection{2.3  Kirkpatrick's Four-Level Evaluation Framework}

Kirkpatrick's [7] four-level training evaluation model, first published in 1959 and refined over the subsequent six decades, provides a hierarchical framework for evaluating the impact of interventions: Level 1 (Reaction) assesses immediate participant response, Level 2 (Learning) measures knowledge or skill acquisition, Level 3 (Behavior) evaluates whether learning transfers to on-the-job behavior change, and Level 4 (Results) measures organizational outcomes [7, 8]. The model has been applied across business, government, military, and healthcare contexts, and remains the most widely used training evaluation framework globally.

We adopt the Kirkpatrick framework as a structural template for evaluating APIP interventions. When a prompt-level, retrieval-level, or model-level intervention is applied to a degraded agent, the check-in evaluation should assess impact at multiple levels: Did the intervention change the agent's immediate output patterns (Level 1/2)? Did it produce durable behavioral change in how the agent handles the failure-mode class (Level 3)? Did it improve the business metrics specified in the behavioral contract---CSAT, task completion, escalation rate (Level 4)? An intervention that improves Level 1/2 metrics without producing Level 3/4 improvement has not remediated the agent; it has merely suppressed the symptom. This hierarchical evaluation discipline is essential for distinguishing genuine recovery from surface-level metric manipulation.

